% ============================================================
%  MTH636 -- Game Theory | Lecture 5
%  Auctions (Cont.) & Max-Min Strategies (Cont.)
% ============================================================
\documentclass[12pt, a4paper]{article}

% ============================================================
%  MTH636 -- Game Theory
%  Notes Preamble — Color-Coded, TikZ-Enhanced
% ============================================================

% --- Core packages ---
\usepackage[T1]{fontenc}
\usepackage[utf8]{inputenc}
\usepackage{lmodern}
\usepackage[margin=2.5cm, top=2.8cm, bottom=2.8cm]{geometry}
\usepackage{microtype}

% --- Math ---
\usepackage{amsmath, amssymb, amsthm, mathtools}
\usepackage{bm}        % bold math
\usepackage{cancel}    % strike-through in math

% --- Color & Boxes ---
\usepackage[dvipsnames, svgnames, x11names, table]{xcolor}
\usepackage[most]{tcolorbox}
\tcbuselibrary{breakable, skins, theorems}

% --- TikZ & Diagrams ---
\usepackage{tikz}
\usepackage{tikz-qtree}
\usetikzlibrary{arrows.meta, positioning, shapes.geometric, fit, backgrounds,
                decorations.pathreplacing, calc, matrix, trees}
\usepackage{pgfplots}
\pgfplotsset{compat=1.18}

% --- Tables ---
\usepackage{booktabs}
\usepackage{multirow}
\usepackage{array}
\usepackage{tabularx}
\usepackage{makecell}

% --- Typography & Layout ---
\usepackage{enumitem}
\usepackage{parskip}
\usepackage{fancyhdr}
\usepackage{titlesec}
\usepackage{titletoc}
\usepackage{hyperref}
\usepackage{graphicx}
\usepackage{float}
\usepackage{caption}

% --- Color Palette ---
\definecolor{defblue}{HTML}{1A5276}      % Definition header
\definecolor{defbluebg}{HTML}{D6EAF8}    % Definition background
\definecolor{thmgreen}{HTML}{1E8449}     % Theorem header
\definecolor{thmbg}{HTML}{D5F5E3}        % Theorem background
\definecolor{exorange}{HTML}{B7770D}     % Example header
\definecolor{exbg}{HTML}{FEF9E7}         % Example background
\definecolor{warnred}{HTML}{922B21}      % Warning/Key Insight header
\definecolor{warnbg}{HTML}{FADBD8}       % Warning background
\definecolor{noticegray}{HTML}{2C3E50}   % Note header
\definecolor{noticebg}{HTML}{EBF5FB}     % Note background
\definecolor{taketawayyellow}{HTML}{7D6608} % Key Takeaway header
\definecolor{takeawaybg}{HTML}{FDFEFE}   % Takeaway background
\definecolor{accentpurple}{HTML}{6C3483} % Accent / proof
\definecolor{lightgray}{HTML}{F2F3F4}    % shading

% ============================================================
%  TCOLORBOX ENVIRONMENTS
% ============================================================

% -- Definition (Blue) --
\newtcolorbox{defbox}[2][]{
  enhanced, breakable,
  colback=defbluebg, colframe=defblue, coltitle=white,
  fonttitle=\bfseries\sffamily,
  title={Definition: #2},
  attach boxed title to top left={yshift=-2mm, xshift=4mm},
  boxed title style={colback=defblue, rounded corners},
  arc=4pt, boxrule=1.2pt,
  left=6pt, right=6pt, top=8pt, bottom=6pt,
  #1
}

% -- Theorem / Proposition (Green) --
\newtcolorbox{thmbox}[2][]{
  enhanced, breakable,
  colback=thmbg, colframe=thmgreen, coltitle=white,
  fonttitle=\bfseries\sffamily,
  title={Theorem: #2},
  attach boxed title to top left={yshift=-2mm, xshift=4mm},
  boxed title style={colback=thmgreen, rounded corners},
  arc=4pt, boxrule=1.2pt,
  left=6pt, right=6pt, top=8pt, bottom=6pt,
  #1
}

% -- Proposition (Green variant) --
\newtcolorbox{propbox}[2][]{
  enhanced, breakable,
  colback=thmbg, colframe=thmgreen!70!black, coltitle=white,
  fonttitle=\bfseries\sffamily,
  title={Proposition: #2},
  attach boxed title to top left={yshift=-2mm, xshift=4mm},
  boxed title style={colback=thmgreen!70!black, rounded corners},
  arc=4pt, boxrule=1.2pt,
  left=6pt, right=6pt, top=8pt, bottom=6pt,
  #1
}

% -- Example (Orange/Yellow) --
\newtcolorbox{exbox}[2][]{
  enhanced, breakable,
  colback=exbg, colframe=exorange, coltitle=white,
  fonttitle=\bfseries\sffamily,
  title={Example: #2},
  attach boxed title to top left={yshift=-2mm, xshift=4mm},
  boxed title style={colback=exorange, rounded corners},
  arc=4pt, boxrule=1.2pt,
  left=6pt, right=6pt, top=8pt, bottom=6pt,
  #1
}

% -- Key Insight / Warning (Red) --
\newtcolorbox{warnbox}[2][]{
  enhanced, breakable,
  colback=warnbg, colframe=warnred, coltitle=white,
  fonttitle=\bfseries\sffamily,
  title={Key Insight: #2},
  attach boxed title to top left={yshift=-2mm, xshift=4mm},
  boxed title style={colback=warnred, rounded corners},
  arc=4pt, boxrule=1.2pt,
  left=6pt, right=6pt, top=8pt, bottom=6pt,
  #1
}

% -- Note (Gray-Blue) --
\newtcolorbox{notebox}[1][]{
  enhanced, breakable,
  colback=noticebg, colframe=noticegray!60, coltitle=white,
  fonttitle=\bfseries\sffamily,
  title={Note},
  attach boxed title to top left={yshift=-2mm, xshift=4mm},
  boxed title style={colback=noticegray!60, rounded corners},
  arc=4pt, boxrule=0.8pt,
  left=6pt, right=6pt, top=8pt, bottom=6pt,
  #1
}

% -- Key Takeaway --
\newtcolorbox{takebox}[1][]{
  enhanced, breakable,
  colback=takeawaybg, colframe=taketawayyellow, coltitle=white,
  fonttitle=\bfseries\sffamily,
  title={Key Takeaway},
  attach boxed title to top left={yshift=-2mm, xshift=4mm},
  boxed title style={colback=taketawayyellow, rounded corners},
  arc=4pt, boxrule=1.2pt,
  left=6pt, right=6pt, top=8pt, bottom=6pt,
  #1
}

% -- Intuition --
\newtcolorbox{intubox}[1][]{
  enhanced, breakable,
  colback=accentpurple!8, colframe=accentpurple!60, coltitle=white,
  fonttitle=\bfseries\sffamily,
  title={Intuition},
  attach boxed title to top left={yshift=-2mm, xshift=4mm},
  boxed title style={colback=accentpurple!80, rounded corners},
  arc=4pt, boxrule=0.8pt,
  left=6pt, right=6pt, top=8pt, bottom=6pt,
  #1
}

% ============================================================
%  PAYOFF MATRIX HELPER
%  Usage: inside tabular, use \payoff{a}{b} for (a,b)
% ============================================================
\newcommand{\payoff}[2]{$(#1,\;#2)$}
\newcommand{\nash}[2]{\underline{$#1$},\underline{$#2$}}  % underline NE payoffs
\newcommand{\Nset}{\mathbb{N}}
\newcommand{\Rset}{\mathbb{R}}
\newcommand{\Sset}{\mathcal{S}}
\newcommand{\Ui}{u_i}

% ============================================================
%  SECTION FORMATTING
% ============================================================
\titleformat{\section}
  {\large\bfseries\sffamily\color{defblue}}
  {\thesection}{1em}{}[\titlerule]

\titleformat{\subsection}
  {\normalsize\bfseries\sffamily\color{thmgreen!80!black}}
  {\thesubsection}{1em}{}

\titleformat{\subsubsection}
  {\normalsize\itshape\sffamily\color{exorange!80!black}}
  {\thesubsubsection}{1em}{}

% ============================================================
%  HEADER / FOOTER
% ============================================================
\pagestyle{fancy}
\fancyhf{}
\fancyhead[L]{\sffamily\small\color{defblue} MTH636 -- Game Theory}
\fancyhead[R]{\sffamily\small\color{defblue} \leftmark}
\fancyfoot[C]{\sffamily\small\thepage}
\renewcommand{\headrulewidth}{0.4pt}

% ============================================================
%  HYPERREF SETUP
% ============================================================
\hypersetup{
  colorlinks=true,
  linkcolor=defblue,
  urlcolor=thmgreen,
  citecolor=accentpurple,
  bookmarksnumbered=true,
  pdftitle={MTH636 Game Theory Notes},
  pdfauthor={MTH636 Student Notes}
}

% ============================================================
%  MISC HELPERS
% ============================================================
\newcommand{\BR}{\text{BR}}        % Best Response
\newcommand{\NE}{\text{NE}}        % Nash Equilibrium
\newcommand{\SPNE}{\text{SPNE}}    % Subgame Perfect NE
\newcommand{\Dom}{\text{dom}}      % Dominance
\newcommand{\mixs}[1]{\sigma_{#1}} % Mixed strategy
\newcommand{\pures}[1]{s_{#1}}     % Pure strategy


% ============================================================
\begin{document}

% Title
\begin{center}
  {\LARGE\bfseries\sffamily\color{defblue} MTH636 -- Game Theory}\\[6pt]
  {\Large\sffamily Lecture 5: Auctions (Cont.) \& Max-Min Strategies}\\[4pt]
\end{center}
\vspace{0.5cm}
\hrule height 1.5pt
\vspace{0.3cm}

% ============================================================
\section{Auctions (Continued)}

Recall from the previous class regarding the \textbf{Second-Price Sealed-Bid Auction}:
\[
  u_i(v_i, b_{-i}) \geq u_i(b'_i, b_{-i}) \quad \forall b_{-i} \in [0, \infty)^{n-1}, \; \forall b'_i \in [0, \infty)
\]
This states that bidding your true valuation $v_i$ weakly dominates any other bid $b'_i$.

\begin{exbox}{3-Player Auction Example}
  Suppose there are three players. For Player 1, their true valuation is $v_1 = 11$.
  Assume the other bids are $b_2 = 10$ and $b_3 = 8$.
  \begin{itemize}
      \item If Player 1 bids $b_1 = 11$ (truthful), they win and pay $\max(10, 8) = 10$. Their payoff is $11 - 10 = 1$.
      \item If they bid lower, say $b_1 = 9$, they lose to Player 2 (who bid 10) and get a payoff of $0$.
      \item If they bid higher, say $b_1 = 15$, they still win and still pay 10. Payoff is still $1$.
  \end{itemize}
  Thus, deviating from the true valuation either doesn't change the payoff or strictly hurts the player.
\end{exbox}

\begin{warnbox}{First-Price \& Third-Price Auctions}
  \begin{itemize}
      \item The above theorem \textbf{does not hold in a First-Price Auction}. In a first-price auction, if you bid your true valuation and win, your payoff is $v_i - v_i = 0$. You are always incentivized to bid lower than your valuation (shade your bid) to obtain a positive surplus.
      \item \textbf{Exercise:} Assume $n \geq 3$ and check if the theorem holds for a \textbf{Third-Price Auction} (where the winner pays the third-highest bid).
  \end{itemize}
\end{warnbox}

% ============================================================
\section{Max-Min Value and Max-Min Strategies (Continued)}

A player plays a strategy $s_i$. The worst-case payoff they could get is:
\[
  \min_{s_{-i} \in S_{-i}} u_i(s_i, s_{-i})
\]
For player $i$, by playing a specific strategy, they can ensure they get at least:
\[
  \max_{s_i \in S_i} \left( \min_{s_{-i} \in S_{-i}} u_i(s_i, s_{-i}) \right)
\]
This value is called the \textbf{max-min value} for player $i$, often denoted $\underline{v}_i$.

And any strategy $s_i^*$ that achieves this, i.e., 
\[
  s_i^* \in \arg\max_{s_i \in S_i} \left( \min_{s_{-i} \in S_{-i}} u_i(s_i, s_{-i}) \right)
\]
is called a \textbf{max-min strategy} for player $i$. If a player plays a max-min strategy, he/she will get at least his max-min value.

\subsection{Infinite Strategy Spaces}
If the strategy spaces are infinite, we use \textit{supremum} and \textit{infimum} instead of max and min. The max-min value is defined as:
\[
  \sup_{s_i \in S_i} \inf_{s_{-i} \in S_{-i}} u_i(s_i, s_{-i})
\]

\begin{exbox}{Guessing Game / Average Game}
  Players submit numbers $x_1, x_2, x_3$. The target $T = T(x_1, x_2, x_3)$ is some function (e.g., the average). 
  The payoff $u_i(x_1, x_2, x_3)$ is $0$ if someone else is closer to the target, and $\frac{1}{|\text{winners}|}$ if there is a tie for being closest.
  
  What is Player 1's max-min value?
  \begin{itemize}
      \item If $x_1 = 10$, the opponents can always choose $x_2, x_3$ such that Player 1 loses. So $\inf_{x_2, x_3} u_1(10, x_2, x_3) = 0$.
      \item If $x_1 = 0$, similarly $\inf_{x_2, x_3} u_1(0, x_2, x_3) = 0$.
  \end{itemize}
  The max-min value is $0$. Therefore, \textbf{any} strategy is a max-min strategy, because no strategy can guarantee a payoff strictly greater than $0$.
\end{exbox}

\begin{notebox}
  For the \textbf{second-price auction}, the max-min value is $0$. Bidding your valuation guarantees at least 0 (you never pay more than your valuation). Playing anything less than your valuation is also a max-min strategy (since it also prevents you from losing money, guaranteeing at least 0).
\end{notebox}

% ============================================================
\section{Max-Min Values and Dominated Strategies}

Suppose Player $i$'s max-min value is $\underline{v}_i$ in game $G$. Let $G'$ be a new game obtained by removing a dominated strategy.

\begin{thmbox}{Effect of Eliminating Strategies}
  \begin{enumerate}
      \item If we remove a strategy $t_i$ of \textbf{Player $i$} that is dominated by some other strategy, Player $i$'s max-min value in $G'$ is \textbf{still exactly} $\underline{v}_i$. (Removing a bad option for yourself doesn't hurt your guaranteed payoff).
      \item If we remove a strategy $t_j$ of \textbf{Player $j$} (an opponent) that is dominated, Player $i$'s max-min value in $G'$ \textbf{may increase} (or stay the same, but it will not decrease). (Removing options from your opponent restricts how much they can hurt you).
  \end{enumerate}
\end{thmbox}

% ============================================================
\section{Introduction to Two-Player Zero-Sum Games}

\begin{defbox}{Two-Player Zero-Sum Game}
  A game $G = \langle N, (S_i)_{i \in N}, (u_i)_{i \in N} \rangle$ is called a two-player zero-sum game if:
  \begin{enumerate}
      \item $N = \{1, 2\}$
      \item $u_1(s_1, s_2) + u_2(s_1, s_2) = 0 \quad \forall (s_1, s_2) \in S_1 \times S_2$
  \end{enumerate}
\end{defbox}

This means whatever Player 1 wins, Player 2 loses. $u_2(s_1, s_2) = -u_1(s_1, s_2)$.
Because of this, we often only write the payoffs for Player 1 in the matrix.

\begin{center}
\renewcommand{\arraystretch}{1.5}
\begin{tabular}{|c|c|c|}
  \hline
  $P_1 \setminus P_2$ & $s_{2,1}$ & $s_{2,2}$ \\
  \hline
  $s_{1,1}$ & $u_1(s_{1,1}, s_{2,1})$ & $u_1(s_{1,1}, s_{2,2})$ \\
  \hline
  $s_{1,2}$ & $u_1(s_{1,2}, s_{2,1})$ & $u_1(s_{1,2}, s_{2,2})$ \\
  \hline
\end{tabular}
\end{center}

\end{document}